\documentclass[12pt,a4paper]{article}
\usepackage[utf8]{inputenc}
\usepackage[T1]{fontenc}
\usepackage{lmodern}
\usepackage[french]{babel}
\usepackage[margin=2.2cm]{geometry}
\usepackage{amsmath,amssymb}
\usepackage{bm}
\usepackage{booktabs}
\usepackage{enumitem}
\usepackage[dvipsnames]{xcolor}
\usepackage[most]{tcolorbox}
\setlength{\parskip}{0.5em}\setlength{\parindent}{0pt}
\setlist{nosep}
\pagestyle{empty}

\newcommand{\Var}{\operatorname{Var}}\newcommand{\Cov}{\operatorname{Cov}}\newcommand{\E}{\mathbb{E}}
\newtcolorbox{intuition}[1][]{colback=Goldenrod!8,colframe=Goldenrod!70!black,fonttitle=\bfseries,title={Intuition},#1}
\newtcolorbox{retenir}[1][]{colback=ForestGreen!6,colframe=ForestGreen!70!black,fonttitle=\bfseries,#1}

\begin{document}

\begin{center}
{\LARGE\bfseries Chapitre 7 -- Sections 6.3 et 6.4}\\[4pt]
{\large Régression avec erreurs AR(1) et solution par les MCG}
\end{center}

\textbf{Rappel : c'est quoi AR(1) ?} L'erreur d'aujourd'hui dépend de celle d'hier :
\[
\epsilon_t=\rho\,\epsilon_{t-1}+u_t,\qquad |\rho|<1,
\]
où $u_t$ est un choc nouveau, imprévisible. Exemple : une grève fait chuter les ventes en janvier ; février sera probablement encore un peu en dessous. $\rho$ mesure à quel point le choc d'hier « déteint » sur aujourd'hui.

% =====================================================================
\section*{6.3 -- Que deviennent les MCO ?}

Modèle : $y_t=x_t'\beta+\epsilon_t$, avec $\epsilon_t=\rho\epsilon_{t-1}+u_t$.

\subsection*{1. Les MCO restent sans biais}
L'absence de biais n'exige que $\E[\epsilon\mid X]=0$ (erreurs non liées aux $x$). Que les erreurs soient corrélées \emph{entre elles} ne change rien.
\begin{intuition}
Les MCO ne visent pas systématiquement trop haut ou trop bas : les erreurs se « suivent », mais elles restent centrées sur zéro.
\end{intuition}

\subsection*{2. Les MCO sont inefficaces}
Les MCO traitent chaque observation comme une information nouvelle et indépendante. Or, avec $\rho>0$, l'observation d'aujourd'hui répète en partie celle d'hier.
\begin{intuition}
C'est comme interroger 100 personnes dont la moitié répètent ce qu'a dit leur voisin : on croit avoir 100 avis, mais on en a moins. Les MCO ne tiennent pas compte de cette redondance ; les MCG, si.
\end{intuition}

\subsection*{3. La variance estimée est fausse, donc les tests aussi}
Le logiciel calcule $\sigma^2(X'X)^{-1}$, une formule qui suppose des erreurs non corrélées : elle \textbf{oublie les covariances}.

\textbf{Exemple} (moyenne de deux observations, exercice 2 du chapitre 7) :
\begin{itemize}
\item formule MCO : $\Var=\sigma_\epsilon^2/2$ ;
\item vraie variance : $\Var=\sigma_\epsilon^2(1+\rho)/2$.
\end{itemize}
Avec $\rho=0{,}5$, la vraie variance est \textbf{1,5 fois plus grande} que celle affichée.
\begin{intuition}
Quand $\rho>0$, les erreurs vont dans le même sens plusieurs périodes de suite : elles ne se compensent pas bien, et l'estimation est plus incertaine que ce que croit le logiciel. Écarts-types \textbf{trop petits}, $t$ \textbf{trop grands}, et on rejette $H_0$ trop souvent.
\end{intuition}

% =====================================================================
\section*{6.4 -- La solution : les MCG (avec $\rho$ connu)}

\textbf{Objectif :} transformer le modèle pour que ses erreurs ne soient plus corrélées, puis appliquer les MCO.

\subsection*{Étape 1 : écrire le modèle à deux dates}
\begin{align*}
\text{aujourd'hui :}\quad & y_t=x_t'\beta+\epsilon_t \tag{38}\\
\text{hier, multiplié par }\rho\text{ :}\quad & \rho y_{t-1}=\rho x_{t-1}'\beta+\rho\epsilon_{t-1} \tag{39}
\end{align*}

\subsection*{Étape 2 : soustraire (39) de (38)}
\[
y_t-\rho y_{t-1}=(x_t-\rho x_{t-1})'\beta+\underbrace{(\epsilon_t-\rho\epsilon_{t-1})}_{=\,u_t}
\]
Donc $\ y_t^*=x_t^{*\prime}\beta+u_t$, avec $y_t^*=y_t-\rho y_{t-1}$ et $x_t^*=x_t-\rho x_{t-1}$ (la \textbf{quasi-différence}).
\begin{intuition}
L'erreur $\epsilon_t$ contient deux morceaux : ce qui est \textbf{hérité} d'hier ($\rho\epsilon_{t-1}$) et le \textbf{choc nouveau} ($u_t$). En soustrayant $\rho$ fois l'équation d'hier, on retire exactement la partie héritée. Il ne reste que $u_t$, sans lien avec le passé : on peut appliquer les MCO. Le $\beta$ estimé est le même ; seules les erreurs ont été « nettoyées ».
\end{intuition}

\subsection*{Étape 3 : le problème de la première observation}
Pour $t=1$, il faudrait $y_1-\rho y_0$, mais $y_0$ n'existe pas. Si on supprime cette observation (\textbf{Cochrane-Orcutt}), il reste $T-1$ observations : on \textbf{perd de l'information}, ce qui compte surtout en petit échantillon.

\subsection*{Étape 4 : garder la première observation (Prais-Winsten)}
On ne peut pas la garder telle quelle, car son erreur n'a pas la même variance que les $u_t$ :
\[
\Var(u_t)=\sigma_u^2,\qquad \Var(\epsilon_1)=\frac{\sigma_u^2}{1-\rho^2}\quad(\text{plus grande}).
\]
\begin{intuition}
$u_t$ est un seul choc. $\epsilon_1$, lui, accumule \textbf{tous les chocs passés} ($u_1+\rho u_0+\rho^2u_{-1}+\dots$), donc il varie davantage.
\end{intuition}
On multiplie la première observation par $\sqrt{1-\rho^2}$ :
\[
\Var\big(\sqrt{1-\rho^2}\,\epsilon_1\big)=(1-\rho^2)\cdot\frac{\sigma_u^2}{1-\rho^2}=\sigma_u^2\ \checkmark
\]
d'où la première équation transformée :
\[
\sqrt{1-\rho^2}\,y_1=\sqrt{1-\rho^2}\,x_1'\beta+\sqrt{1-\rho^2}\,\epsilon_1 .
\]
Sa variance est maintenant la même que celle des autres observations.

\subsection*{Le modèle complet transformé (44)}
\begin{center}
\renewcommand{\arraystretch}{1.4}
\begin{tabular}{@{}lcc@{}}
\toprule
& $t=1$ & $t>1$ \\
\midrule
$y^*$ & $\sqrt{1-\rho^2}\,y_1$ & $y_t-\rho y_{t-1}$ \\
$x^*$ & $\sqrt{1-\rho^2}\,x_1$ & $x_t-\rho x_{t-1}$ \\
erreur & $\sqrt{1-\rho^2}\,\epsilon_1$ & $u_t$ \\
\bottomrule
\end{tabular}
\end{center}
On applique les MCO à $y^*$ sur $x^*$ : ce sont les \textbf{MCG}, l'estimateur efficace.

\textbf{Piège pratique :} la constante est elle aussi transformée. Sa colonne de 1 devient $\sqrt{1-\rho^2}$ pour $t=1$ et $1-\rho$ pour $t>1$. On ne rajoute donc pas de constante habituelle : on régresse sur cette colonne transformée.

\textbf{Limite :} tout cela suppose $\rho$ \textbf{connu}. En pratique, on l'estime (section 6.5).

% =====================================================================
\section*{Ce qu'il faut retenir}

\begin{retenir}[title={6.3 -- Les MCO avec des erreurs AR(1)}]
\begin{enumerate}
\item \textbf{Sans biais} : les erreurs restent non liées aux $x$.
\item \textbf{Inefficaces} : les observations se répètent en partie, et les MCO ne le savent pas.
\item \textbf{Tests faux} : $\sigma^2(X'X)^{-1}$ oublie les covariances ; avec $\rho>0$, écarts-types trop petits et rejets de $H_0$ trop fréquents.
\end{enumerate}
\end{retenir}

\begin{retenir}[title={6.4 -- La solution (MCG, $\rho$ connu)}]
\begin{enumerate}
\item \textbf{Quasi-différence} : $y_t-\rho y_{t-1}=(x_t-\rho x_{t-1})'\beta+u_t$. On enlève la partie héritée ; il reste le choc nouveau $u_t$.
\item \textbf{Première observation} : multipliée par $\sqrt{1-\rho^2}$ pour avoir la même variance (Prais-Winsten). La supprimer, c'est Cochrane-Orcutt (perte d'information).
\item \textbf{Ensuite} : MCO sur le modèle transformé $=$ MCG, l'estimateur efficace.
\item \textbf{Limite} : il faut connaître $\rho$ ; en pratique on l'estime (section 6.5).
\end{enumerate}
\end{retenir}

\end{document}
